\documentclass{article}
\usepackage[UTF8]{ctex}
\usepackage{amsmath,amssymb,amsfonts}
\usepackage{enumitem}
\usepackage{geometry}
\geometry{a4paper, margin=1in}
% === 错题本排版优化（不参与题目标识） ===
\setlength{\parindent}{0pt}
\setlength{\parskip}{0.35em plus 0.10em minus 0.08em}
\setlength{\emergencystretch}{3em}
\setlist{itemsep=0.25em, topsep=0.40em, parsep=0pt}
\allowdisplaybreaks[2]
% === 错题本排版优化结束 ===
\begin{document}
\section*{南京师范大学2025年数学分析真题}
\section*{92．南京师范大学2025年数学分析试卷}

\subsection{题目1}
1． \(\lim _{x \rightarrow 1}\left(\frac{1}{x-1}-\frac{1}{\ln x}\right)\) ．

\subsection{题目2}
2． \(\lim _{n \rightarrow \infty} n\left(\frac{1}{n^{2}+1}+\frac{1}{n^{2}+2^{2}}+\cdots+\frac{1}{n^{2}+n^{2}}\right)\) ．

\subsection{题目3}
3． \(\lim _{n \rightarrow \infty} \frac{n+1}{\sqrt[n]{n!}}\) ．

\subsection{题目4}
4．设 \(f(x)=\int_{0}^{x} \sin (x-t) e^{-t} \mathrm{~d} t\) ，求 \(f^{\prime}(x)\) ．

\subsection{题目5}
5． \(\int e^{x} \cos ^{2} x \mathrm{~d} x\).

\subsection{题目6}
6．讨论 \(\int_{0}^{+\infty} \frac{x^{m}}{1+x^{n}} \mathrm{~d} x\) 的敛散性，其中 \(m, n \in \mathbb{N}_{+}\)。

\subsection{题目7}
7．设函数 \(f\) 具有二阶连续偏导数，令 \(z=f(x+y, \arctan (x y))\) ，求 \(\frac{\partial^{2} z}{\partial x \partial y}\) ．

\subsection{题目8}
8．计算曲面积分 \(\iint_{\sum}(x+y+z)^{2} \mathrm{~d} S\) ，其中 \(\sum: x^{2}+y^{2}+z^{2}=R^{2}(R>0)\) ．

\subsection{题目9}
1．设 \(a>0, b>0, a_{1}=a, a_{n+1}=\frac{1}{2}\left(a_{n}+\frac{b}{a_{n}}\right), n \in \mathbb{N}_{+}\)．证明：数列 \(\left\{a_{n}\right\}\) 收敛并计算 \(\lim _{n \rightarrow \infty} a_{n}\) ．

\subsection{题目10}
2．设函数 \(f(x)\) 在 \([a, b]\) 上连续，在 \((a, b)\) 上可导，且 \(f(a)=f(b)=0\) ．证明：存在 \(\xi \in(a, b)\) ，使得

\[
f^{\prime}(\xi)-2 f(\xi)=0 .
\]

三．（15 分）证明函数 \(y=\sqrt{x}\) 在 \([0,+\infty)\) 上一致连续．

四．（15 分）求函数项级数 \(\sum_{n=1}^{\infty} \frac{n x}{(1+x)(1+2 x) \cdots(1+n x)}(x \geq 0)\) 的和函数，并讨论其在 \([1,+\infty)\) 上的一致收敛性。

五．（15 分）设 \(f(u)\) 和 \(g^{\prime}(v)\) 在 \(\mathbb{R}\) 上连续，\(L: x^{2}+y^{2}=4\) ，取逆时针．\(D: x^{2}+y^{2} \leq 4\) ，且

\[
\iint_{D}(x+y) g^{\prime}(x-y) \mathrm{d} x \mathrm{~d} y=1
\]

计算 \(I=\int_{L}\left[f\left(x^{2}+y^{2}\right)+g(x-y)\right](x \mathrm{~d} x+y \mathrm{~d} y)\) ．

六．（15 分）计算

\[
I=\iint_{\sum} 2 x^{3} \mathrm{~d} y \mathrm{~d} z+2 y^{3} \mathrm{~d} z \mathrm{~d} x+3\left(z^{2}-1\right) \mathrm{d} x \mathrm{~d} y
\]

其中 \(\sum\) 是曲面 \(z=1-x^{2}-y^{2}(z \geq 0)\) 的上侧。

七．（15 分）设 \(x_{n} \neq 0, n \in \mathbb{N}_{+}\)，证明：若 \(\lim _{n \rightarrow \infty} n\left(1-\frac{\left|x_{n+1}\right|}{\left|x_{n}\right|}\right)=a>1\) ，则 \(\sum_{n=1}^{\infty} x_{n}\) 绝对收敛．

八．（15 分）解答如下问题。

\begin{itemize}
  \item[（1）] 叙述实数系完备性基本定理中的致密性定理和 Cauchy 收敛准则．
  \item[（2）] 利用致密性定理证明 Cauchy 收敛准则的充分性部分。
\end{itemize}
\end{document}
